% ============================================
% 通用数模论文 LaTeX 导言区
% 适用于大部分中文数学建模竞赛
% ============================================

% === 文档类与中文支持 ===
\documentclass[12pt,a4paper]{article}
\usepackage[UTF8]{ctex}
\usepackage{fontspec}

% === 页面布局 ===
\usepackage[top=2.5cm,bottom=2.5cm,left=2.8cm,right=2.8cm]{geometry}
\usepackage{fancyhdr}
\pagestyle{fancy}
\fancyhf{}
\fancyfoot[C]{\thepage}
\renewcommand{\headrulewidth}{0.4pt}

% === 数学宏包 ===
\usepackage{amsmath,amssymb,mathtools}
\usepackage{bm}
\usepackage{nicefrac}

% === 图表 ===
\usepackage{graphicx}
\usepackage{float}
\usepackage{booktabs}
\usepackage{multirow}
\usepackage{longtable}
\usepackage{subcaption}

% === 算法 ===
\usepackage[ruled,linesnumbered]{algorithm2e}

% === 列表 ===
\usepackage{enumitem}

% === 颜色与超链接 ===
\usepackage{xcolor}
\usepackage[colorlinks=true,linkcolor=black,citecolor=blue,urlcolor=blue]{hyperref}

% === 参考文献 ===
\usepackage{gbt7714}
\bibliographystyle{gbt7714-numerical}

% === 代码高亮 ===
\usepackage{listings}
\lstset{
  basicstyle=\small\ttfamily,
  keywordstyle=\color{blue}\bfseries,
  commentstyle=\color{gray},
  stringstyle=\color{red},
  numbers=left,
  numberstyle=\tiny\color{gray},
  frame=single,
  breaklines=true,
  captionpos=b,
  tabsize=4
}

% === 自定义命令 ===
\newcommand{\vect}[1]{\bm{#1}}
\newcommand{\mat}[1]{\mathbf{#1}}
\newcommand{\diff}{\mathrm{d}}
\newcommand{\pd}[2]{\frac{\partial #1}{\partial #2}}
\newcommand{\dd}[2]{\frac{\diff #1}{\diff #2}}
